\documentclass[11pt]{article}

\usepackage{fontspec}
\setmainfont{TeX Gyre Pagella}

\usepackage[a4paper,margin=1.15in]{geometry}
\usepackage{microtype}
\usepackage[english]{babel}

\usepackage{amsmath}
\usepackage{amssymb}
\usepackage{amsthm}

\theoremstyle{plain}
\newtheorem{theorem}{Theorem}
\newtheorem{definition}{Definition}
\newtheorem{proposition}{Proposition}
\newtheorem{conjecture}{Conjecture}
\newtheorem{example}{Example}
\newtheorem{corollary}{Corollary}

\usepackage[dvipsnames]{xcolor}
\usepackage[colorlinks=true,linkcolor=violet,citecolor=violet,urlcolor=violet]{hyperref}
\usepackage{csquotes}

\setlength{\parindent}{1.2em}
\setlength{\parskip}{0pt}

\title{\textbf{The Compression Fallacy}\\[0.3em]
\large Why Shorter Is Not the Same as Better Understood}
\author{Flyxion}
\date{August 2026}

\begin{document}

\maketitle

\begin{abstract}
\noindent
A widely held assumption, rarely stated but frequently acted on, treats the shortest description of a pattern as the correct measure of how well that pattern is understood: a good model is a compressed model, and compression is the currency of understanding. This paper argues that the assumption outruns what has actually been established about the relationship between description length and understanding, and proves a non-identification result stating exactly how far: for a fixed object, its description length under one ontology and its distinguishing capacity are independently realizable to any target values, and no ordering of description lengths across two ontologies constrains the ordering of their distinguishing capacities. Description length, in other words, does not identify distinguishing capacity, without further assumptions supplied from outside the comparison. A previously proposed result, the Deficit Proxy Theorem, is examined as a candidate additional-assumption case: it claims that under a faithful-encoding condition, perfect compression guarantees perfect distinguishing capacity, which would pin down one point on the otherwise unidentified surface---but closer inspection shows that even this special case is not fully secured by the argument given for it, since it treats generating an object cheaply and distinguishing that object within a specified domain as though they were the same achievement, and closing the gap between them would need a further bridge condition the source result does not supply. The Forth programming language supplies a worked case in which description length and distinguishing capacity come apart cleanly: Forth code is short not because it approaches the Kolmogorov-optimal description of the task it solves, but because it is reachable, at low cost, from an already-accumulated library of admissible operations---a distinct and prior achievement the brevity of the resulting code does not, by itself, report, and the paper's separation of a representation's stored magnitude from the machinery it presupposes generalizes this observation past Forth to trained models and to structure-preserving cryptographic constructions alike. The paper closes by proposing that a representation's distinguishing capacity, tested directly and, where possible, under deliberate intervention on the domain, is the quantity that should be measured wherever it and description length can be pulled apart.
\end{abstract}

\newpage
\tableofcontents
\bigskip

\section{An Old Definition, Asked of a New Case}

The computer scientist Daniel Hillis once offered a definition of the information content of a pattern: the amount of information in a pattern of bits equals the length of the smallest computer program capable of generating those bits. The definition is, in substance, Kolmogorov complexity restated for a general audience, and Hillis's own statement of it claims nothing beyond that restatement. What this paper is concerned with is a further claim the definition has often been taken, by others, to support, though it is not one Hillis himself is on record as drawing: that a system which produces a shorter description of something has thereby understood that thing better than a system producing a longer one. The claim is rarely defended directly. It is simply assumed, in the background of arguments about which model of a phenomenon is the better model, which representation of a domain is the more insightful representation, and---increasingly, in discussions of large trained models---which network has learned the most genuine structure rather than having merely memorized surface statistics.

Forth, an unusually terse programming language whose devotees have long pointed to Hillis's definition as evidence of the language's virtue, supplies an unusually clean test case for the assumption, because Forth's brevity is well documented, uncontroversial, and produced by a mechanism that is easy to state precisely: word-factored Forth code is short because a Forth programmer builds a library of named operations---words---and subsequent code invokes them rather than re-deriving them (Ziembik 2025). A task solved in Forth after such a library exists can be several times shorter, in raw character count, than the same task solved in a general-purpose language without access to an equivalent library. The question this paper asks is not whether Forth code is short; that is not in dispute. The question is what that shortness is evidence of, and whether it is evidence of the thing Hillis's definition, read the way it is usually read, would have it be evidence of.

This paper argues that it is not, and that the same gap separating Forth's brevity from Kolmogorov-optimal description separates, more consequentially, model compactness from model understanding wherever the comparison is drawn. Section~\ref{sec:apparatus} introduces the formal apparatus needed to state the gap precisely, developed independently of any concern with Forth or with machine learning. Sections~\ref{sec:lossless}--\ref{sec:relocation} clear three preliminary confusions the apparatus invites---lossless compression's relationship to distinguishing capacity, the reference machinery any description length is measured against, and where complexity goes when a description shrinks without that machinery being counted. Section~\ref{sec:nonident} then proves this paper's mathematical center, a non-identification result stating that description length does not determine distinguishing capacity anywhere except at one identifiable point, and Section~\ref{sec:proxy} recovers the previously available Deficit Proxy Theorem as exactly that point, purchased by an additional assumption rather than assumed as the paper's foundation. Section~\ref{sec:forth} works through the Forth case as a concrete instance of the gap. Section~\ref{sec:ml} turns to the contemporary target: the informal but pervasive assumption, in discussions of trained models, that shorter or sparser representations are better understood ones. Section~\ref{sec:converse} takes up a converse case from cryptography, in which a field explicitly preserves an algebra of operations rather than raw representation size. Section~\ref{sec:overcompression} names the point at which compression becomes irreversible identification, and its relationship to sufficiency relative to a specific question. Section~\ref{sec:curricula} argues that engineered, deliberately simplified learning environments and productive detours through greater complexity during problem-solving are further evidence against treating compression progress as monotonic with understanding. Section~\ref{sec:intervention} proposes testing distinguishing capacity under deliberate intervention rather than only at rest. Section~\ref{sec:conclusion} proposes what should be measured instead, where the two quantities can be pulled apart and are not.

\section{Coding Length and Distinguishing Capacity}
\label{sec:apparatus}

A general framework for measuring representational capacity (Flyxion 2026)\footnote{This paper departs from the practice, followed elsewhere in this series, of not citing the author's own prior work, for the same reason given in a companion paper: the argument here depends on a specific prior formal result, and describing that result without attribution would make it impossible to check against its source.} takes as its basic object not a set of items to be described, nor a description language considered on its own, but the pairing of the two together with a notion of when two items count as the same for the purposes at hand. Formally, an \emph{ontological triple} $(X,\sim,T)$ consists of a set $X$, an equivalence relation $\sim$ on $X$ specifying which elements of $X$ are treated as observationally indistinguishable, and an ontology $T$: a description language, or template hierarchy, determining which descriptions of elements of $X$ are actually reachable, as opposed to merely existing in principle. The distinction between what a description language can express in principle and what it can reach in practice, given the templates and operations it has already accumulated, is the hinge the rest of this paper turns on.

Two deficits are defined over this structure, and they measure different things. The \emph{coding deficit} of an object $x$, written $\delta_T(x)$, is the excess length of the shortest description of $x$ reachable within $T$ over the Kolmogorov-optimal description length of $x$---the length of the absolute shortest program capable of generating $x$, considered independently of any particular ontology. The \emph{distinguishability deficit} of $x$, written $\Delta_T(x)$, is a different quantity entirely, built from two further pieces of notation worth naming on their own: let $D_T(x)$ denote the \emph{distinguishing capacity} of $T$ at $x$, the number of other objects that can be separated from $x$ by descriptions reachable in $T$, and let $D^*(x)$ denote the maximum such capacity attainable by any ontology whatsoever. The distinguishability deficit is then $\Delta_T(x) = D^*(x) - D_T(x)$: the gap between what $T$ actually distinguishes at $x$ and what the best possible ontology could. The coding deficit asks how much longer a description is than it needs to be. The distinguishability deficit asks how much distinguishing power has been lost. Nothing in the definitions forces these two quantities to move together, and the relationship between them, rather than being assumed, is the object a further theorem is required to establish.

Two further pieces of structure need to be made explicit before proceeding, since the constructions below depend on keeping them visibly separate. First, because $\sim$ already specifies which elements of $X$ count as observationally equivalent for the purposes of the problem, it is useful to pass to the \emph{semantic quotient} $Q = X/\!\sim$ and treat an ontology's encoder as acting on classes $q \in Q$ rather than on arbitrary representatives of $X$. An ontology $T$ decomposes as a triple $T = (E_T, M_T, \mathcal{O}_T)$: an encoder $E_T : Q \to \{0,1\}^*$ giving each semantic class its reachable description; background machinery $M_T$, the accumulated templates, interpreter, or structure constants required to produce or use those descriptions in the first place; and an \emph{operational repertoire} $\mathcal{O}_T$, the set of operations or relations computable directly over the encoded representations $\{E_T(q) : q \in Q\}$ without first decoding back to $Q$. Nothing in this decomposition is new relative to the definition already given; it only names three pieces of it separately, and $\mathcal{O}_T$ in particular will not receive a full formal theory in what follows---it is introduced because Section~\ref{sec:converse} needs it named, not because this paper builds machinery for it beyond that use. $L_T(q) = |E_T(q)|$ recovers the length notion already introduced. A second length, $L(T)$, measures the cost of $M_T$ itself in some fixed reference language, and the total $L_{\mathrm{total}}(q) = L(T) + L_T(q)$ will matter once Section~\ref{sec:relative} asks what $L_T(q)$ is being measured relative to.

The encoder also induces a second equivalence relation, this one on $Q$ rather than on $X$: $q \sim_T r \iff E_T(q) = E_T(r)$, recording which semantically distinct classes the ontology nevertheless collapses. Write $\mathcal{P}_T = Q/\!\sim_T$ for the resulting partition, and $[q]_{\sim_T}$ for $q$'s block within it. The pointwise distinguishing capacity, restated against this notation, is $D_T(q) = |Q| - |[q]_{\sim_T}|$: the number of semantic classes $T$ actually separates $q$ from. Keeping $\sim$ and $\sim_T$ visibly distinct matters throughout what follows: $\sim$ specifies which distinctions matter, prior to any ontology; $\sim_T$ records which of those distinctions a particular ontology's encoder actually preserves.

It is worth collecting the three quantities an ontology now carries into a single definition, since the rest of the paper is organized around keeping them separate rather than letting any one substitute for another.

\begin{definition}[Three coordinates of a representation]
\label{def:threecoord}
For an ontology $T = (E_T, M_T, \mathcal{O}_T)$, distinguish three properties: the \emph{description magnitude} $L_T(q) = |E_T(q)|$; the \emph{distinguishing structure} $\mathcal{P}_T = Q/\!\sim_T$; and the \emph{operational repertoire} $\mathcal{O}_T$, the operations reachable using $M_T$ directly over encoded representations. These measure different properties of $T$. Equality or improvement in one does not, without an additional coupling assumption supplied from outside the definitions, imply equality or improvement in either of the others: writing $\not\Rightarrow$ for ``does not imply,''
\[
L_T \not\Rightarrow \mathcal{P}_T, \qquad L_T \not\Rightarrow \mathcal{O}_T, \qquad \mathcal{P}_T \not\Rightarrow \mathcal{O}_T,
\]
stated as three separate non-implications rather than a chain, since a chained arrow would suggest only adjacent pairs are being denied.
\end{definition}

Proposition~\ref{prop:nonident}, proved in Section~\ref{sec:nonident}, is the formal content behind the first of these three: for any $\ell_1 < \ell_2$, none of
\[
D_{T_1}(q) > D_{T_2}(q), \qquad D_{T_1}(q) = D_{T_2}(q), \qquad D_{T_1}(q) < D_{T_2}(q)
\]
is ruled out by $L_{T_1}(q) < L_{T_2}(q)$ alone. Section~\ref{sec:converse} supplies the case for the third, where Silva's construction and ordinary CKKS can agree completely on $L_T$ and $D_T$ while disagreeing on $\mathcal{O}_T$.

\section{Lossless Compression Does Not Imply Ontological Economy}
\label{sec:lossless}

Before turning to the machine-learning case, it is worth addressing the objection this apparatus invites immediately: doesn't lossless compression preserve all the information in a representation, so that a shorter lossless encoding must preserve every distinction the longer one did? At the level of static distinguishability this case is in fact straightforward, and it is worth being precise about why rather than reaching for a different notion of capacity to explain the intuition the objection is trying to express. A ZIP archive can become dramatically shorter than the file it compresses while remaining perfectly invertible: its encoder is injective on $Q$, so $\sim_T$ collapses nothing, every block of $\mathcal{P}_T$ is a singleton, and $D_T(q) = |Q|-1$ for every $q$---maximal, by the definitions already given. Lossless compression is accordingly not a counterexample to static distinguishability preservation, and nothing in Proposition~\ref{prop:nonident} claims otherwise.

What the ZIP case exposes instead is a third quantity this paper has not yet needed: $\mathcal{O}_T$, the operational repertoire introduced in Section~\ref{sec:apparatus}. That two encoded objects remain distinct, in the static sense $D_T$ measures, does not imply that every relation of interest between their \emph{uncompressed contents} remains directly computable over the encoded representations themselves. Two compressed files can be told apart trivially---their bitstrings differ---without any geometric or relational property of what they contain being available for comparison until both are decompressed; if decompression is not itself among the operations $\mathcal{O}_T$ makes native, then those relations sit outside $\mathcal{O}_T$ even though $D_T$ is already maximal. The relevant separation is accordingly three-way rather than two-way: description magnitude ($L_T$), distinguishing structure ($\mathcal{P}_T$, or $D_T$ pointwise), and operational repertoire ($\mathcal{O}_T$) are logically independent properties of an ontology, and a lossless archive is the case in which the first two are already exemplary while the third remains, until an inverse transform is paid for, essentially empty.

The converse case rules out the opposite inference and returns to familiar territory. An extremely verbose database---one that stores every field of every record with no compression whatsoever, in a schema an order of magnitude larger than the information it strictly requires---can distinguish every member of its domain perfectly despite being terribly compressed: $\delta_T(x)$ large, $\Delta_T(x) = 0$. Nobody encountering such a database concludes that its size is evidence against its usefulness, because nobody was using size as a proxy for distinguishing power in the first place; the intuition that a bloated schema is not thereby a poorly understood one is already the intuition this paper is asking to be extended to the cases---sparse latent codes, monosemantic feature dictionaries---where the opposite inference, size taken as evidence of understanding, is currently made routinely. Section~\ref{sec:nonident} states, in full generality, why neither inference is licensed by description length alone.

\section{Compression Relative to What?}
\label{sec:relative}

A further assumption is buried in any claim that a representation is ``short'': short relative to what reference machinery. Kolmogorov complexity has its invariance theorem, which guarantees that the choice of universal machine affects $L^*(x)$ by at most an additive constant independent of $x$---but that guarantee is asymptotic, and practical comparisons between representational systems are made nowhere near the regime in which an additive constant can be safely ignored. Writing $L_{\mathrm{visible}}(q) = L_T(q)$ for the quantity ordinarily reported and $L_{\mathrm{system}}(q) = L(M_T) + L_T(q)$ for the full cost, the point is that a comparison of the first carries no guarantee about the second:
\[
L_{T_2}(q) < L_{T_1}(q) \quad \not\Rightarrow \quad L(M_{T_2}) + L_{T_2}(q) < L(M_{T_1}) + L_{T_1}(q).
\]
Forth already makes the point vivid without needing to be reformulated: the token \texttt{FOO} is three characters, but the definition of \texttt{FOO} together with the words it depends on can represent thousands of characters of accumulated machinery. Reporting $L_{\mathrm{visible}}(q)$---the three characters---without reporting $L(M_T)$---the cost of the machinery \texttt{FOO} presupposes---is not measuring a smaller quantity than $L_{\mathrm{system}}(q)$; it is measuring a different, incomplete one, and the difference between the two is exactly the accumulated word library Section~\ref{sec:forth} examines directly. The same separation generalizes past Forth without alteration: a short latent code in a trained autoencoder can depend on a decoder with billions of parameters, and reporting the latent code's length while omitting the decoder's is the identical move under a different name.

\section{Where Did the Complexity Go?}
\label{sec:relocation}

The pattern in Section~\ref{sec:relative} generalizes to a lemma worth stating on its own, because it clarifies what a short $L_T(x)$ is actually evidence of.

\begin{proposition}[Machinery Displacement]
Let $Q = X/\!\sim$ be finite. There exists an ontology $T_2 = (E_2, M_2)$ in which every $q \in Q$ has a fixed-width description of length $\lceil \log_2 |Q| \rceil$, while the machinery required to interpret those descriptions contains the corresponding lookup structure. Consequently, reducing the per-object description length does not by itself establish that representational complexity has been eliminated rather than displaced into background machinery.
\end{proposition}

\begin{proof}
Enumerate $Q = \{q_0, \dots, q_{n-1}\}$, $n = |Q|$, and define $E_2(q_i) = \mathrm{bin}_{\lceil \log_2 n \rceil}(i)$. Let $M_2$ contain the correspondence between these indices and the classes they denote. Every encoded object then has the same description length $\lceil \log_2 n \rceil$, while interpreting those descriptions presupposes the correspondence stored in $M_2$. The cost visible in $L_{T_2}(q)$ can thus be reduced by moving representational work into $M_2$, with no conservation law for total description length required for the conclusion: different coding schemes and interpreters can change $L_{\mathrm{total}}$ in either direction, and nothing here claims otherwise.
\end{proof}

Compression magnitude measured without $L(T)$, in other words, cannot by itself distinguish whether structure has been eliminated from a description or merely relocated into machinery that is not being counted; $L_T(q)$ alone cannot tell which happened. The question worth asking of any short description is accordingly not \emph{how short is $d(x)$?} but \emph{what machinery must already exist for $d(x)$ to recover $x$?} A shell command invokes an entire operating system. A Forth word invokes its dictionary. A latent vector invokes a decoder. A theorem invokes its definitions and lemmas. An encrypted Clifford geometric product, taken up in Section~\ref{sec:converse}, invokes a pre-established table of structure constants. In every one of these cases, a description length reported in isolation, without the machinery it presupposes, is not a smaller number than the full cost; it is the wrong number, with the missing part relocated rather than eliminated.

\section{Description-Length Non-Identification}
\label{sec:nonident}

The preceding three sections motivate, without yet formally establishing, the claim this section proves: that description length alone, considered apart from the machinery and partition structure that produced it, does not determine distinguishing capacity, in either direction and at no point short of the single extreme the Deficit Proxy Theorem addresses. What follows is a genuine non-identification result, not a counterexample to one proposed correspondence: it establishes that fixing $L_T(q)$ leaves $D_T(q)$ completely free, and that no ordering of description lengths across two ontologies constrains the ordering of their distinguishing capacities.

\begin{proposition}[Description-Length Non-Identification]
\label{prop:nonident}
Let $Q = X/\!\sim$ be finite with $|Q| \ge 3$, and let $q \in Q$. Then:
\begin{enumerate}
\item[(i)] For every pair $(\ell, k)$ with $\ell \ge 1$ and $0 \le k \le |Q| - 1$, there exists an ontology $T$ with $L_T(q) = \ell$ and $D_T(q) = k$.
\item[(ii)] For each relation $R \in \{<, =, >\}$, there exist $T_1, T_2$ with $L_{T_1}(q) < L_{T_2}(q)$ and $D_{T_1}(q) \; R \; D_{T_2}(q)$.
\end{enumerate}
\end{proposition}

\begin{proof}
Let $n = |Q|$. For (i), choose a subset $C \subseteq Q$ containing $q$ with $|C| = n - k$. Define the encoder so that every element of $C$ receives the same description $s$ of length $\ell$. Assign every element of $Q \setminus C$ a description distinct from $s$ and from every other such description; these descriptions may have arbitrary lengths. Since $\{0,1\}^*$ is infinite, such an assignment always exists. The encoder-induced class of $q$ is therefore exactly $C$, so $D_T(q) = |Q| - |C| = n - (n-k) = k$, and $|E_T(q)| = \ell$, so $T$ realizes the requested pair $(\ell,k)$. Statement (ii) follows by applying (i) twice, with $\ell_1 < \ell_2$ and $k_1, k_2$ chosen so that $k_1 \; R \; k_2$.
\end{proof}

A fully worked instance makes the construction checkable without appeal to the general argument, and, per Proposition~\ref{prop:nonident}, is stated pointwise throughout: fix $Q = \{a,b,c\}$ and compare $L_T(a)$ against $D_T(a)$ across five encoders.

\begin{center}
\begin{tabular}{lllll}
Encoder & $E(a)$ & $E(b)$ & $E(c)$ & $(L_T(a),\,D_T(a))$ \\
\hline
$F_1$ & 000 & 000 & 000 & $(3,\,0)$ \quad total collapse \\
$F_2$ & 0 & 1 & 1 & $(1,\,2)$ \quad $a$ alone \\
$F_3$ & 00 & 00 & 01 & $(2,\,1)$ \quad $a$ grouped with $b$ \\
$F_4$ & 00 & 01 & 10 & $(2,\,2)$ \quad injective \\
$F_5$ & 0000 & 0001 & 0010 & $(4,\,2)$ \quad injective, padded \\
\end{tabular}
\end{center}

Each row fixes $q=a$ and varies only which classes $a$ is grouped with, exactly as the proof's construction does. $F_3$ against $F_4$: equal length ($\ell=2$), unequal capacity ($k=1$ vs.\ $k=2$). $F_2$ against $F_1$: $\ell=1<3$ with $k=2>0$, shorter length and strictly higher capacity. $F_3$ against $F_5$: $\ell=2<4$ with $k=1<2$, the case ordinary usage means by ``compression loses information,'' present here as one populated cell among several rather than the only possibility. $F_4$ against $F_5$: $\ell=2<4$ with $k=2=2$, a shorter injective encoding against a longer injective one, identical capacity. Every cell of the ordering table is realized by an explicit construction at the same fixed point $a$; none is left as an existence claim or drawn from a shifting statistic.

Proposition~\ref{prop:nonident} says exactly that $L_T(q)$, the \emph{compression magnitude}, and $\mathcal{P}_T = Q/\!\sim_T$, the \emph{compression topology}, are independent coordinates of $T$:

\begin{center}
\fbox{\parbox{0.85\linewidth}{\centering compression magnitude does not determine compression topology}}
\end{center}

fixing one leaves the other free to take any admissible value. Setting $k = |Q|-1$ (so $\Delta_T(q)=0$) and letting $\ell \to \infty$ recovers Appendix~\ref{app:counterexamples}'s padding example as a single point on a now fully characterized surface, rather than an isolated existence claim standing on its own.

The one thing doing real work in the proof is that $T$ is unconstrained: no cost links a class's description length to its size. This is not an assumption smuggled in for the purpose of the proof; it is precisely the freedom already granted by Section~\ref{sec:apparatus}'s definition of an ontology, and precisely what real systems exhibit---padding, hash truncation, and a bloated-but-short-on-average library all occupy different corners of the $(\ell,k)$ plane this proposition describes. Faithful encoding, the assumption the source Deficit Proxy Theorem adds, removes exactly one corner of that freedom: it rules out $\ell$ at its minimum, $L^*(q)$, coexisting with $k$ below its maximum---an optimal-length code that still fails to distinguish much, the representational analogue of a hash collision. That is the sense in which $\delta_T(q)=0$ would force $\Delta_T(q)=0$ under that assumption, examined more carefully, and found wanting even at that single point, in Section~\ref{sec:proxy}. What matters here is only that faithful encoding, whatever it secures, secures nothing about any $\ell > L^*(q)$, which is the entire range in which real representational systems, including every example this paper discusses, actually operate. Proposition~\ref{prop:nonident} holds everywhere, unconditionally, with no dependence on faithful encoding or on any other assumption about how $T$ encodes.

A further consequence follows immediately once a sequence of ontologies, rather than a single one, is in view---the formal target the discussion of compression progress in Section~\ref{sec:curricula} needs.

\begin{proposition}[Compression Progress Non-Identification]
\label{prop:dynamic}
Let $(T_t)_{t=0}^{m}$ be a sequence of ontologies over a fixed finite $Q$, and $q \in Q$. Write $\Delta L_t = L_{T_{t+1}}(q) - L_{T_t}(q)$ and $\Delta D_t = D_{T_{t+1}}(q) - D_{T_t}(q)$. Without additional constraints linking description length to encoder-induced partitions,
\[
\operatorname{sgn}(\Delta L_t) \; \not\Rightarrow \; \operatorname{sgn}(\Delta D_t).
\]
In particular, $\Delta L_t < 0$ is compatible with each of $\Delta D_t < 0$, $\Delta D_t = 0$, and $\Delta D_t > 0$.
\end{proposition}

\begin{proof}
Apply Proposition~\ref{prop:nonident} independently at times $t$ and $t+1$, choosing $\ell_{t+1} < \ell_t$ while choosing $k_t, k_{t+1}$ via part (i) to realize any desired ordering between them.
\end{proof}

Differentiating an unidentified relationship with respect to time does not identify it. A compression-progress account that reads $\Delta L_t < 0$ as evidence of improving understanding, or a temporary $\Delta L_t > 0$ as evidence of regress, is committing Proposition~\ref{prop:nonident}'s error once per time step, not avoiding it by moving to a derivative; Section~\ref{sec:curricula} uses exactly this consequence against the strongest version of that account.

\section{The Proxy Theorem as a Conditional Result}
\label{sec:proxy}

A previously proposed result, the Deficit Proxy Theorem, is worth examining here as a candidate additional-assumption case of Proposition~\ref{prop:nonident}---the kind of result that pins down one point on an otherwise unidentified surface by adding a specific, checkable assumption. Stated in its source, it says that under an assumption of faithful encoding---that every distinction a description language can express costs at least one bit to express---the distinguishability deficit is bounded above by a constant multiple of the coding deficit: $\Delta_T(x) \le C\,\delta_T(x)$, with the special case $\delta_T(x) = 0 \implies \Delta_T(x) = 0$ claimed as a direct consequence: an ontology achieving the Kolmogorov-optimal description of an object would thereby also achieve the maximum possible distinguishing capacity for it.

Closer inspection, carried out in full in Appendix~\ref{app:proxy-theorem}, shows that even this special case is not secured by the argument given for it. The step from Kolmogorov-optimality to the claimed conclusion treats generating $x$ cheaply and distinguishing $x$ from an ambient domain as though they were the same achievement; they are not, and closing the gap between them would require an additional condition tying $L^*(x)$ to a maximally distinguishing encoding specifically, a condition the source result does not supply and this paper does not attempt to supply either. The theorem is presented here, accordingly, as a previously proposed conditional result under examination rather than as an established point this paper's own argument depends on: even its zero-deficit implication requires a further bridge condition connecting Kolmogorov-optimal generation of an object to maximal distinction of that object within a specified domain, and that bridge has not, on the evidence available, been built.

The converse---that a distinguishability deficit of zero forces a coding deficit of zero, or more generally that reducing an already-small coding deficit further continues to track improvements in distinguishing capacity---fares no better; the source theorem states it only as a conjecture, resting on the further assumption that optimal coding uses every distinction available to it, with no independent argument offered for that assumption and no obvious reason to expect it holds for representational systems built by processes other than deliberate, distinction-preserving design. Proposition~\ref{prop:nonident} shows why no version of this tracking should be expected in general: description length and distinguishing capacity vary independently across their entire joint range, and nothing examined in this section identifies even one point on that range without a gap of its own.

The asymmetry matters because the popular reading of Hillis's definition uses the coding deficit as a graded, continuous stand-in for the distinguishability deficit across its whole range, treating even the strongest available version of the Deficit Proxy Theorem as though it licensed that continuous use. It would not license it even if its own special case were secure: an inference from perfect compression to perfect distinguishing capacity is not an inference that a further reduction in an already-imperfect coding deficit constitutes evidence of a further reduction in distinguishing capacity, still less that the amount of the one measures the amount of the other. Section~\ref{sec:forth} argues that the unconstrained region Proposition~\ref{prop:nonident} describes is, in fact, close to the ordinary condition of a mature, well-factored code library.

\section{Forth and Reachable Complexity}
\label{sec:forth}

The source framework's own introduction, developed independently of this paper's concerns, names but does not develop a third notion relevant here: \emph{reachable complexity}, the shortest description of an object expressible from a description language's current template hierarchy, which may be strictly longer than the Kolmogorov-optimal description---longer, that is, than the true value of $L^*(x)$ against which the coding deficit is measured. Forth supplies reachable complexity's first concrete instance. A Forth programmer accumulates, over the life of a project, a library of named words: short, composable operations built from more primitive ones, each capturing an operational distinction the programmer has found useful to keep separately addressable. Code written after such a library exists is short because it is reachable, cheaply, from that library---not because it approaches $L^*(x)$, the description length that would be achieved by an ontology with no library at all, starting from nothing, optimizing purely for brevity against the raw task.

This is not a minor qualification. It inverts the naive reading of Hillis's definition as applied to Forth. The naive reading takes Forth's brevity as evidence that Forth programs approach the information-theoretic floor of the tasks they solve, and hence as evidence that Forth, among languages, comes closest to expressing the true content of a computation. The reachable-complexity reading takes the same brevity as evidence of something else: that a well-factored library of admissible operations was available to be reused. Two Forth programmers solving the same task with different accumulated word libraries will produce differently short code for it, and the difference in length reports which library was available, not which programmer better understood the task. A companion result, developed independently of any concern with Forth, states the point in general form: no valid compression mapping can be constructed without prior identification of the invariant structure of the space being compressed, and that identification is itself achieved only through a prior process of expansion---exploration, iteration, and the discovery of which structure is invariant in the first place (Flyxion 2026b). Forth's terseness is compression in exactly this downstream sense: it presupposes the expansion that built the word library, and reports that expansion's prior existence rather than reporting anything about the Kolmogorov-optimal floor of the task the short code happens to solve.

None of this counts against Forth, and it is worth being explicit that it does not, since the point of this section is not to relitigate a decades-old argument about programming-language design. A large, well-factored word library achieving low $\delta_T(x)$ relative to itself, while $L^*(x)$ remains a separate and largely inaccessible quantity, is exactly what a mature, high-distinguishing-capacity representational system should look like. The coding deficit measured relative to the library is doing real work in this case; what it cannot do, without the further, unproved direction of the Proxy Theorem, is stand in for a claim about how close the resulting code sits to the absolute informational floor of the problem, or, more importantly for what follows, for a claim about how much of the task's true distinguishing structure the library actually preserves.

\section{The Same Move in Machine-Learned Representations}
\label{sec:ml}

The informal assumption this paper has been arguing against is not confined to programming-language advocacy; it is closer to background theory in contemporary discussions of trained models, where a shorter, sparser, or more compressed internal representation is routinely treated as evidence of deeper or more genuine understanding, and description-length reduction has been a recognized modeling virtue since well before large neural networks existed, formalized in the minimum-description-length principle (Rissanen 1978). The assumption's most explicit and most influential statement extends compression progress well past a technical modeling virtue and treats it as the underlying currency of understanding itself: a unifying account of curiosity, creativity, and aesthetic interest across infants, mathematicians, artists, and artificial systems alike, in which subjective beauty and interestingness are identified with the rate at which an observer's compression of its data is currently improving (Schmidhuber 2008). On this account, the drive to compress better is not merely a useful modeling heuristic but the mechanism generating science, art, music, and humor alike, which makes it the clearest available target for the argument of this paper: if the relationship between coding deficit and distinguishability deficit is only proven in one direction and only at the extreme, then a theory that identifies the felt sense of understanding, or beauty, or interest with compression progress at every point along the curve, rather than only at the point where compression is complete, is claiming a great deal more than the formal relationship between the two quantities actually licenses. Sparse-coding and disentanglement objectives are frequently motivated by an appeal to compression as a proxy for the discovery of a domain's true generative factors, and post-hoc interpretability work often treats the discovery of a smaller, cleaner set of directions or features within a model's activations as evidence that those features constitute the model's genuine ontology, on grounds close to Hillis's, and close to Schmidhuber's: the shorter description is taken to be the truer one, and the steeper the recent gain in shortness, the more interesting or better understood the underlying structure is taken to be. Two examples make the pattern concrete rather than merely characterized. The beta-VAE architecture increases a single compression pressure---a weighting term on the model's latent-code cost, pushed higher than in an ordinary variational autoencoder---specifically in order to produce latent representations described as more interpretable and as more faithfully capturing a domain's independent underlying factors of variation, with the increase in compression pressure offered as the mechanism by which the improvement in interpretability is obtained (Higgins et al.\ 2017). Sparse-autoencoder interpretability work on language models motivates the move to sparse, low-dimensional dictionaries of features in comparable terms: individual neurons are described as polysemantic and therefore poorly understood, while the sparse, more compressed feature directions recovered by the autoencoder are described as monosemantic and correspondingly more interpretable, with the compression step treated as what makes the improved interpretability available rather than as an independent, separately justified engineering choice (Bricken et al.\ 2023). In neither case is the inference from compression to understanding stated as an unproven assumption requiring a caveat; in both, it is the stated motivation for the method.

The apparatus of Sections~\ref{sec:apparatus} and~\ref{sec:proxy} applies to this case without modification, and the case exposes a further difficulty the Forth case does not. It is worth being precise about what is being targeted before applying it. Many researchers working on sparse representations, mechanistic interpretability, and minimum-description-length methods already treat compression as a useful heuristic rather than as a theorem about understanding, and nothing in this paper is addressed to a practitioner who already holds the distinction Section~\ref{sec:proxy} makes formal. The target is narrower and more common than any claim about what a field collectively believes: compression, or a reduction in a representation's apparent size, is routinely offered \emph{as evidence for} deeper or more genuine understanding, in papers, talks, and informal argument alike, without the caveat that only one direction of the relevant relationship, and only its extreme case, has actually been established. It is the inference, not the community, that this paper is aimed at, and the inference survives even where individual practitioners would, if asked directly, decline to endorse it.

The Deficit Proxy Theorem's special case, examined in Section~\ref{sec:proxy} and Appendix~\ref{app:proxy-theorem}, depends on a faithful-encoding assumption: that every distinction a representation expresses costs at least one unit of the representation's capacity to express. Distributed and superposed representations, of the kind large trained networks are now widely understood to use---in which a single direction in activation space participates in encoding more than one feature, and a single feature is encoded across more than one direction---are precisely representations for which faithful encoding, in the sense the assumption requires, need not hold. Where it fails, even the weakened, already-incompletely-established special case no longer applies, and a reduction in a representation's measured or apparent description length carries no guaranteed relationship to its distinguishing capacity at all, in either direction. The inference from compression to understanding is, in this setting, not merely unsupported past some extreme; it is not obviously licensed anywhere along the range where trained models actually operate. The same displacement occurs in text summarization systems whose short introductory sentences function as pointers into a larger source: the visible sentence carries little of the recoverable content itself, while the machinery that selects the supplementary material carries the mapping from the pointer to what it makes reachable.

None of this is an argument that compression is irrelevant to understanding, and the paper does not make that argument. A system with a smaller coding deficit is, all else equal, a more efficient one, and efficiency is its own legitimate goal. The argument is narrower and, for that reason, more defensible: that the graded, continuous use of description length as a running proxy for understanding, throughout the actual working range of real representational systems, is not supported by the result usually cited in its favor---and, per Section~\ref{sec:proxy}, is not securely supported even at the single extreme that result is usually taken to establish---and in the specific case of distributed neural representations may not be supported by any comparable result at all.

\section{Preserving Structure Rather Than Recovering It}
\label{sec:converse}

The examples above show fields treating compression as evidence of understanding. A recent proposal in homomorphic encryption offers a different kind of case: not a field measuring the same quantity this paper has been discussing, but a field independently discovering that a representation's stored values do not by themselves determine what can be done with it---which turns out to be a sharper instance of this paper's thesis than a simple compression-versus-size contrast would be.

Silva's Clifford fully homomorphic encryption (Clifford FHE) targets a familiar practice in applied cryptography: geometric data---vectors, orientations, rotations, volumes---is ordinarily flattened into independent scalar components before encryption, since standard homomorphic schemes operate on scalars (Silva 2026). It is tempting to read this as a direct instance of $\delta_T$-reduction, with the reported $250{,}000\times$ expansion between a 64-byte plaintext multivector and its roughly 16-megabyte encrypted form standing in for the coding-deficit gap this paper's apparatus measures. That reading does not hold up. The expansion Silva reports is ciphertext expansion: redundancy that ring-learning-with-errors encryption introduces for cryptographic security, not a comparison of description lengths in the sense $\delta_T(x) = L_T(x) - L^*(x)$ is defined here. A large ciphertext is not evidence of a large coding deficit, and the size figures should not be read as one.

The more interesting distinction survives this correction, but it is not one this paper's static apparatus, $L_T$ and $D_T$, was built to measure---it is an instance of the third coordinate introduced in Section~\ref{sec:apparatus}, $\mathcal{O}_T$. Silva's construction still encrypts each of a multivector's eight components separately---eight CKKS ciphertexts, exactly as an ordinary flattened encoding would produce, and by the definitions given here $D_T$ is unaffected by this choice either way: the component values are equally present, equally distinct, under both schemes. What differs is $\mathcal{O}_T$: a fixed table of structure constants specifying how the eight encrypted components combine under the geometric product, so that rotation, projection, and the wedge and inner products are native operations over the ciphertexts, computable directly rather than only after decryption. Two encodings can hold identical component values, and identical static distinguishing capacity, and still differ entirely in whether a consumer of those values can compute the geometric relations between them without first decrypting. ``Discards'' overstates what ordinary flattened CKKS does to those relations: nothing about scalar CKKS makes a rotation or an orientation mathematically unrecoverable from the eight ciphertexts, since the component values are all still there and could in principle be recombined by hand using scalar homomorphic operations. What ordinary CKKS lacks is the Clifford structure as a \emph{native, first-class} part of $\mathcal{O}_T$: recovering a rotation from flattened ciphertexts requires re-deriving, term by term, an algebra that Clifford FHE simply supplies as an available operation. Silva's own framing matches this reading directly: the paper states that associativity, distributivity, and rotor conjugation remain available over ciphertexts specifically because the geometric product is implemented as a homomorphism, not because the underlying values were otherwise unreachable.

This is a sharper case than the toy counterexample of Appendix~\ref{app:counterexamples} makes on its own, and it motivates rather than merely illustrates $\mathcal{O}_T$'s presence in this paper's apparatus: two ontologies can agree completely on $L_T$ and on $D_T$---identical stored components, identical static distinguishing capacity---while disagreeing entirely on what can be computed over those components without an inverse transform. $E_{T_1}(q) \approx E_{T_2}(q)$ at the level of stored values is compatible with $\mathcal{O}_{T_1} \ne \mathcal{O}_{T_2}$, and it is the second disagreement, not the first, that Silva's construction is actually about. Silva's structure constants are, in this paper's vocabulary, part of $\mathcal{O}_T$: they are what make the geometric product native to the encrypted representation, and their absence from ordinary CKKS is what makes recovering rotation and orientation there a matter of re-derivation rather than direct computation, even though the same eight numbers are, in principle, sitting in the ciphertexts either way.

The connection to Section~\ref{sec:forth}'s treatment of Forth is direct rather than merely analogical. Forth's terseness was traced there to an accumulated word library determining what computations are cheaply reachable from the interpreter's current vocabulary, not to any special proximity to the informational floor of the task. Silva's structure constants play the same role for geometric operations on encrypted data: what a representation can be made to do is set by $\mathcal{O}_T$, the operational apparatus carried alongside it, and a comparison of stored sizes or even static distinguishing capacity alone, in either case, misses the part of the representation that actually does the work.

Two qualifications keep the case appropriately narrow. First, the baseline comparison Silva reports---a classical multilayer perceptron on three mean-coordinate features reaching 35\% accuracy and failing under rotation, against 99\% for the Clifford encoding---is not a controlled ablation of flattening alone. The two systems receive different input features: the MLP baseline sees only a centroid, while the multivector encoding is built from richer handcrafted statistics, including average radial distance, principal orientation planes, and a volume measure. The comparison is suggestive of what an operationally impoverished representation costs, but it does not isolate the effect of discarding the geometric algebra from the effect of discarding the additional summary statistics, and should not be cited as a controlled measurement of either alone. Second, the reported 99\% encrypted-versus-100\%-plaintext result uses handcrafted weights on a three-class toy problem, a limitation the paper states outright, with full gradient-based training left as future work. What the case supports is the structural claim---that reachable distinctions depend on preserved operations, not merely preserved values---not the specific accuracy figures generalizing beyond this proof of concept.

\section{Overcompression and Catastrophic Identification}
\label{sec:overcompression}

The apparatus so far has treated $\Delta_T(x)$ and $\delta_T(x)$ as quantities that can diverge without either being pathological on its own. It is worth naming directly what happens at the point where divergence becomes irreversible. Suppose $x_1 \ne x_2$ but an encoding maps both to the same representation, $E(x_1) = E(x_2)$. Compression has crossed an epistemic boundary at that pair: unless information external to $E$ remains available, the distinction between $x_1$ and $x_2$ cannot subsequently be reconstructed from the encoding alone. This is precisely the equivalence relation already in use, $x \sim_E y \iff E(x) = E(y)$, examined at the level of what it asserts rather than merely what it computes. Every compression scheme implicitly asserts that the distinctions collapsed within its induced equivalence classes do not matter, whether or not that assertion was made deliberately. The question worth asking of any compressed representation is accordingly not \emph{how much was compressed} but \emph{which differences the compression declared equivalent}---a question about $\mathcal{P}_T$, the compression topology of Section~\ref{sec:nonident}, not about $L_T(x)$, its magnitude.

This reframing has an immediate, useful qualification, worth stating as its own definition rather than left as a caveat: a collapsed distinction is not automatically a discarded one, if the collapse was made relative to a specific question rather than to the domain as a whole.

\begin{definition}[Task-relative equivalence]
For a task $\tau : X \to Y$---a function returning whatever answer the task actually requires---define $x \sim_\tau y \iff \tau(x) = \tau(y)$. An encoding $E$ \emph{preserves the distinctions $\tau$ requires} when $E(x) = E(y) \implies x \sim_\tau y$, equivalently $\sim_E \, \subseteq \, \sim_\tau$.
\end{definition}

\begin{definition}[Overcompression relative to a task]
$E$ \emph{overcompresses relative to $\tau$} if there exist $x, y \in X$ with $E(x) = E(y)$ but $\tau(x) \ne \tau(y)$: equivalently, $\sim_E \, \not\subseteq \, \sim_\tau$.
\end{definition}

This is sharper than calling every collision ``catastrophic identification,'' since it makes plain that some collisions are exactly what a task requires and others are not, and that the difference is not visible from $E$ alone---it depends on $\tau$. A statistic $S(X)$ can throw away enormous quantities of raw information while retaining everything necessary for inference about a specified parameter $\theta$; if $S$ is sufficient for $\theta$, then $\sim_S \, \subseteq \, \sim_\theta$ by definition, so $S$ does not overcompress relative to $\theta$, however much it overcompresses relative to $X$ itself. Compression, on this reading, is not opposed to understanding; it is meaningful only relative to a prior decision about which distinctions must survive. The equivalence classes $\sim_S$ induced by a sufficient statistic are not evidence of poor understanding of $\theta$---they are exactly what good understanding of $\theta$, and nothing else, looks like. The failure this section is naming is not compression itself but overcompression relative to some further task the representation is silently expected to also serve: a representation with $\sim_E \subseteq \sim_\theta$ mistaken for one with $\sim_E \subseteq \sim_X$.

\section{Curricula Must Widen}
\label{sec:curricula}

A further difficulty for any theory that identifies understanding, or interest, or progress with monotonic compression is that the environments in which learning actually occurs are not found in a state of nature; they are built, and built deliberately simple, for reasons that have nothing to do with the learner's understanding having reached some genuine informational floor. Each scaffold described in this section temporarily sacrifices environmental richness in order to preserve the distinctions a learner can currently acquire: not every distinction the environment could in principle present, but the subset a learner at the current stage can actually absorb, with the remainder deliberately withheld rather than lost. Children are not raised inside the full chaos of an unfiltered natural environment, in which every relevant variable changes at once and at every timescale simultaneously; they are raised inside environments deliberately constrained to present regularities slowly, one or a few at a time, precisely because a learner who cannot yet compress much of anything will fail to learn from an environment offering nothing but unstructured, simultaneous complexity. The same engineering shows up outside pedagogy proper, for reasons that are frankly economic rather than developmental. Standardized building materials---dimensional lumber cut to a small number of fixed sizes, sheet goods produced in uniform panels---constrain the built environment to a narrow vocabulary of regular shapes because mass production rewards uniformity, and the resulting simplicity of built spaces is a byproduct of manufacturing economics, not evidence that simple, rectilinear spaces are closer to the true shape of anything. Roads are built straight and their surfaces held deliberately unchanging across long distances in part because the vehicles traveling them have limited suspension travel and no capacity to adapt their gait to terrain the way a walking animal does; the road's low geometric entropy compensates for the vehicle's narrow tolerance, and says nothing about the terrain the road happens to cross.

In each case, the low entropy of the environment is a scaffold fitted to the current, limited distinguishing capacity of whatever is operating within it, not a report on the true complexity of the domain the environment sits inside. This matters for the compression-progress account directly, because a scaffold is only useful for as long as the learner has not yet outgrown it. Once a learner is no longer acquiring anything new from a fixed training environment, whether that learner is a child who has mastered a curriculum or a system whose compression of its current data has plateaued, further progress requires increasing the entropy of what is presented---widening $D_T(x)$ toward the distinctions the scaffold had been withholding---not continuing to compress what is already compressed as far as it usefully can be. A curriculum that only ever simplifies, and never widens, is a curriculum that terminates in a plateau rather than in understanding. This is difficult to reconcile with a theory that identifies the drive toward understanding with the drive toward compression progress at every point along the learner's trajectory, since the trajectory that theory describes is one of steadily increasing compression, while the trajectory that actually produces continued learning is one of alternating compression and deliberately reintroduced complexity. Proposition~\ref{prop:dynamic} gives this tension a formal target rather than leaving it as an observation about trajectories: a compression-progress account reads $\Delta L_t < 0$ as evidence of improving understanding and a temporary $\Delta L_t > 0$ as evidence of regress, but neither inference is licensed once $\operatorname{sgn}(\Delta L_t) \not\Rightarrow \operatorname{sgn}(\Delta D_t)$ at every single time step---and the productive-detour phenomenon this section describes is exactly a case realizing $\Delta L_t > 0$ with $\Delta D_t > 0$, a temporary expansion accompanying rather than undermining improved distinguishing capacity.

A related and sharper difficulty appears within problem-solving itself, on timescales far shorter than a curriculum. Solving a problem often requires making it more complicated for a while before it can be made simple: introducing an auxiliary variable that has no simplifying role on its own but opens a path to one, embedding a problem in a higher-dimensional space where a solution becomes visible before projecting the result back down, or otherwise moving, temporarily and deliberately, into what might be called a higher-dimensional logic relative to the problem's original statement. A compression-progress account, tracking description length as it evolves along the solver's trajectory, would register this intermediate expansion as a local reversal---a temporary loss of interestingness or beauty, on that account's own terms---precisely at the point where the solver is doing the genuine work that will shortly produce the solution. A theory of understanding that cannot distinguish a productive detour through greater complexity from a genuine loss of progress has mismeasured the trajectory it is trying to describe.

The curriculum-widening argument is not merely anecdotal; it has an independently developed empirical literature behind it that was arrived at for unrelated reasons. Developmental research on the explore--exploit trade-off argues that the distinctive length of human childhood is itself an evolved solution to a tension between exploring broadly and exploiting what is already known, with early cognition characterized by wide, high-entropy hypothesis search that is progressively narrowed only as the demands of goal-directed exploitation increase, and with young learners repeatedly found to explore more broadly than older learners and adults precisely where broad exploration is advantageous (Gopnik 2020). On this account childhood is not an inefficient, not-yet-compressed version of adult cognition working its way toward an adult's compression; it is a distinct cognitive strategy, favoring the preservation of distinctions a purely exploitative learner would discard, for exactly as long as the developmental trade-off favors exploration over efficiency. A learner whose only objective were continuous compression progress would converge on a narrow, exploited hypothesis space too early, at the first local efficiency gain available, and would never acquire the broader distinguishing capacity that later skilled performance depends on. Human development instead allocates an extended period in which exploration is favored over immediate efficiency, which is a second, independent line of evidence for the same structural claim Section~\ref{sec:curricula}'s engineered examples make: that widening a curriculum, not narrowing it further, is what continued distinguishing capacity requires at exactly the point where a monotonic compression-progress account would expect further narrowing.

\section{Understanding Under Intervention}
\label{sec:intervention}

The recommendation this paper has been building toward---measure distinguishing capacity directly rather than inferring it from description length---is still, as stated, a claim about a static quantity, $D_T(x)$, computed once against a fixed domain. A stronger and more operational version is available: rather than asking how many distinctions a representation preserves at rest, ask which distinctions survive when the domain is deliberately perturbed. Write $D_T(x \mid I)$ for the distinguishing capacity of $T$ at $x$ after an intervention $I$---a rotation, a change to an otherwise irrelevant coordinate, a novel composition of familiar primitives, an out-of-distribution combination the representation was not directly trained or built against---and define the \emph{intervention robustness}
\[
R_T(x, I) = D_T(x \mid I) - D_T(x),
\]
the change in distinguishing capacity the intervention induces. Two representations with identical $L_T(x)$, or even identical $D_T(x)$ measured before any intervention, can behave completely differently under $I$: formally,
\[
L_{T_1}(x) = L_{T_2}(x) \quad \text{and} \quad D_{T_1}(x) = D_{T_2}(x) \quad \not\Rightarrow \quad D_{T_1}(x \mid I) = D_{T_2}(x \mid I).
\]
One representation preserves the same relations after the perturbation, the other does not, and the difference is invisible to any measurement taken only at rest---agreement on the first two coordinates of Definition~\ref{def:threecoord} says nothing about behavior under $I$, which is a further, independent fact about $T$.

Silva's construction supplies a motivating case that is unusually clean because the claim being tested is architectural rather than statistical. Rotational equivariance in the geometric neural network is not an empirical regularity discovered by training on many rotated examples; it follows from the algebraic construction itself, since $W \otimes (R v \tilde R) = R(W \otimes v)\tilde R$ holds as an identity of the geometric product, independent of any particular weights---in the notation just introduced, $R_T(x, I) = 0$ for $I$ a rotation is a claim provable from $T$'s algebra rather than measured empirically. A representation whose distinguishing capacity is claimed to be preserved under rotation can accordingly be tested directly, by rotating the input and checking whether the relations the representation is supposed to preserve still hold, rather than inferred from the size of the representation or the accuracy of a single held-out test set. This is a sharper standard than the static $D_T(x)$ this paper has otherwise used, and the sharper standard is the useful one: a representation's distinguishing capacity is revealed by which distinctions survive admissible transformations of the domain, not merely by how many bits are required to encode the domain's resting state.

\section{What Should Be Measured Instead}
\label{sec:conclusion}

Proposition~\ref{prop:nonident} establishes that coding length and distinguishing capacity can come apart across their entire joint range except at one identified point; the Forth and machine-learning cases in Sections~\ref{sec:forth} and~\ref{sec:ml} show the gap is not merely possible but ordinary. A practice that measures the first while reporting conclusions about the second is accordingly measuring the wrong quantity almost everywhere it is applied. The correction this paper proposes is not a rejection of compression as a modeling virtue---a system with a smaller coding deficit is, all else equal, a more efficient one, and efficiency is its own legitimate goal---but a separation of that virtue from claims about understanding, which should be evaluated against the distinguishability deficit directly wherever the two quantities can be independently assessed: does the representation actually preserve the ability to separate the objects, cases, or outcomes that matter for the task, rather than merely occupying fewer bits while doing so.

A useful example of this distinction already appears in evaluation practice for text summarization. Systems designed to produce very short summaries confront the obvious fact that brevity can preserve a coarse overview while eliminating distinctions present in the source. Some summarization work therefore supplements automatic overlap measures such as ROUGE with separate human judgments of properties such as completeness and cohesion, rather than treating the automatic score as a sufficient measurement of summary quality (Sotudeh and Goharian 2022).\footnote{The system's own name for its introductory-sentence pointers is apt enough to borrow without much alteration: ``follow the torches left in the introductory passages, and mine the depths with renewed clarity.''} Structurally, this is the measurement discipline advocated here, and it is worth being precise about what kind of example it is: ROUGE is not a compression metric, and nothing about it measures description length. It is a readily computable proxy observable, cheap where the property of actual interest is not, and the paper's own practice of measuring it separately from human judgments of completeness is an instance of a broader principle this essay's apparatus is one case of: $P(T) \not\Rightarrow Q(T)$ for a cheap proxy $P$ and a target property $Q$, absent some independently established relation between them. Whether the source is length or lexical overlap, the discipline is the same---evaluate the property of interest directly wherever it can be pulled apart from what is convenient to compute.

The same system's central mechanism supplies a second, more literal instance of this paper's specific apparatus. Sotudeh and Goharian's summarizer selects short introductory sentences and uses them as pointers into the rest of the source document, with a trained selection network $M$ mapping each pointer to the supplementary sentences it is taken to indicate. Writing $\mathrm{Reach}(s;M) = \{x \in X : x \text{ is selectable from } s \text{ using } M\}$ for the set a pointer sentence $s$ makes reachable under that machinery, the relationship between a pointer's own length and its reach is exactly the relationship Proposition~\ref{prop:nonident} describes:
\[
|s_1| < |s_2| \quad \not\Rightarrow \quad |\mathrm{Reach}(s_1; M_1)| \lessgtr |\mathrm{Reach}(s_2; M_2)|.
\]
A short pointer sentence can have a large reach given a sufficiently capable $M$; a longer one can have almost none. The visible sentence is $L_T$; the trained selector is $M_T$; what the system is actually engineered to get right is $\mathrm{Reach}$, which the sentence's own length does not report.

This is a harder quantity to measure than description length, which is precisely why description length has been used as its substitute. The Forth case suggests one route past the difficulty: ask not how short a representation is, but what library of prior, admissible distinctions its shortness is reachable from, and whether that library, examined on its own terms, actually covers the distinctions the task requires. The machine-learning case suggests a second: ask whether a representation's apparent compactness survives an audit for superposition and faithful encoding before treating that compactness as evidence of anything beyond itself. Section~\ref{sec:intervention} suggests a third, more demanding than either: perturb the domain and ask which distinctions survive the perturbation, rather than trusting a static count taken at rest. None of the three is as convenient as reading a bit count off a trained model or a line count off a Forth program. Convenience is exactly what the coding deficit offers, and exactly what Proposition~\ref{prop:nonident} shows is not available almost anywhere it is currently taken to be. Appendices~\ref{app:definitions} through~\ref{app:counterexamples} restate the formal apparatus this recommendation rests on, so that the case for measuring distinguishing capacity directly can be checked against exactly what has, and has not, been proved about its relationship to coding length, rather than taken on the paper's word; Appendix~\ref{app:consequences} draws out what the non-identification result implies for representation comparison more generally, once description magnitude, distinguishing structure, and operational repertoire are recognized as independently variable rather than aspects of a single underlying quality.

\newpage
\appendix

\section{Formal Definitions}
\label{app:definitions}

For reference, the definitions used throughout the paper, restated together and without surrounding argument.

An \emph{ontological triple} is $(X,\sim,T)$, where $X$ is a set of objects, $\sim$ is an equivalence relation on $X$ specifying observational indistinguishability, and $T$ is an ontology: a description language or template hierarchy determining which descriptions of elements of $X$ are actually reachable, as distinct from which descriptions merely exist in principle.

A description of $x \in X$ is \emph{reachable within $T$} if it can be constructed using only the templates, operations, or primitives $T$ makes available. The \emph{reachable length} $L_T(x)$ is the length of the shortest description of $x$ reachable within $T$. The \emph{Kolmogorov-optimal length} $L^*(x) = K(x)$ is the length of the absolute shortest program capable of generating $x$, with no restriction to any particular ontology.

The \emph{coding deficit} is $\delta_T(x) = L_T(x) - L^*(x)$, always non-negative since $L^*(x)$ is a minimum over all ontologies.

The \emph{distinguishing capacity} $D_T(x)$ is the number of other objects separable from $x$ by descriptions reachable in $T$. The \emph{maximum distinguishing capacity} $D^*(x)$ is the largest such capacity attainable by any ontology. The \emph{distinguishability deficit} is $\Delta_T(x) = D^*(x) - D_T(x)$, likewise always non-negative.

\emph{Faithful encoding} is the assumption that every distinction a description language can express costs at least one bit of description length to express: no distinction is expressible for free.

Since $\sim$ already specifies which elements of $X$ are treated as identical, it is convenient to work on the \emph{semantic quotient} $Q = X/\!\sim$ directly. An ontology decomposes as a triple $T = (E_T, M_T, \mathcal{O}_T)$: an encoder $E_T : Q \to \{0,1\}^*$, background machinery $M_T$, and an operational repertoire $\mathcal{O}_T$, the operations computable directly over $\{E_T(q):q\in Q\}$ without decoding back to $Q$. The encoder induces $q \sim_T r \iff E_T(q) = E_T(r)$, with $\mathcal{P}_T = Q/\!\sim_T$ the resulting partition and $D_T(q) = |Q| - |[q]_{\sim_T}|$ restated pointwise against it.

\section{The Deficit Proxy Theorem}
\label{app:proxy-theorem}

\begin{theorem}[Deficit Proxy, restated]
Under faithful encoding, $\Delta_T(x) \le C\,\delta_T(x)$ for a system-dependent constant $C > 0$. In particular, $\delta_T(x) = 0 \implies \Delta_T(x) = 0$.
\end{theorem}

\begin{proof}[Proof sketch, as given in the source result]
Under faithful encoding, each distinction inaccessible to $T$ at $x$---each of the $\Delta_T(x)$ separations $T$ fails to make that some ontology achieving $D^*(x)$ would make---would require at least one additional bit in $T$'s description of $x$ to flag, were $T$ to be extended to express it. The number of such extensions needed is bounded by $\Delta_T(x)$, so the excess length $\delta_T(x)$ needed to close the gap is bounded below by $\Delta_T(x)$ up to the constant $C$ scaling the maximum per-distinction cost, giving $\Delta_T(x) \le C\,\delta_T(x)$.

For the special case: if $\delta_T(x) = 0$, then $T$ already achieves the Kolmogorov-optimal description length for $x$. Under faithful encoding, every distinction that would cost a bit to express is therefore already expressible in $T$---an ontology paying the minimum possible length has no unspent length with which additional, unexpressed distinctions could be hiding---so no distinction is inaccessible: $\Delta_T(x) = 0$.
\end{proof}

\emph{A caveat on the special case, worth stating rather than passing over.} The step from Kolmogorov-optimality to ``no unspent length for hidden distinctions'' is not, as written, a direct consequence of $L^*(x)$'s definition. Generating $x$ and distinguishing $x$ from the rest of a specified domain are different tasks: a shortest program for $x$ is not thereby guaranteed to encode every relation between $x$ and every other element of $X$, and $L^*(x) = K(x)$ says nothing by itself about $D^*(x)$. Closing this gap rigorously would need an additional condition tying $L^*(x)$ specifically to a maximally distinguishing encoding, rather than to Kolmogorov-optimality alone---a condition the source result does not supply. This paper does not attempt to supply it either, and now, with Proposition~\ref{prop:nonident} established independently in Section~\ref{sec:nonident}, does not need to: the Deficit Proxy Theorem is presented here as an external result, worth stating because it is the identified point the wider non-identification result is compatible with, not as a result this paper's own argument depends on.

\begin{conjecture}[Full Deficit Correspondence, not established]
Under faithful encoding, $\Delta_T(x) = 0 \iff \delta_T(x) = 0$.
\end{conjecture}

The forward direction is Theorem~1 above, subject to the caveat already noted. The reverse direction --- that a distinguishability deficit of zero forces a coding deficit of zero --- requires the further assumption that optimal coding uses every distinction available to it: that an ontology achieving maximum distinguishing capacity does so via a description that is also length-optimal, rather than via a longer description that happens to preserve every distinction without compressing them efficiently. No argument for this further assumption is offered in the source theorem, and this paper does not supply one either. The direction informal practice actually relies on runs the other way, from $\delta_T(x) \approx 0$ to $\Delta_T(x) \approx 0$: low coding deficit is what is directly observed, in a bit count or a line count, and good distinguishing capacity is what is being inferred from it. That is exactly the direction the caveat above already shows is unsecured even at the theorem's own claimed special case, without yet invoking the reverse direction's further, separately unsupported assumption at all.

\section{Reachable Complexity, Formalized}
\label{app:reachable}

The reachable length introduced in Appendix~\ref{app:definitions} is precisely the formal object informally called reachable complexity in Section~\ref{sec:forth}:
\[
L_T(x) = \min\{\,|d| : d \in T,\ d \mapsto x\,\},
\]
the length of the shortest description $d$, drawn from the templates and operations available in $T$, that generates $x$. By definition of $L^*(x)$ as a minimum over all possible ontologies,
\[
L_T(x) \ge L^*(x) = K(x),
\]
with equality only when $T$ happens to contain a description of $x$ that is itself Kolmogorov-optimal --- a reachability condition on $T$, not a general property of well-factored ontologies. A mature library can drive $L_T(x)$ very close to $L^*(x)$ for the objects it was built around, without this implying anything about $L_{T'}(x)$ for a differently-built library $T'$, or about how close either library sits to $L^*(x)$ for objects outside the range its accumulated templates were shaped by. Forth's brevity, restated in this notation, is a claim about $L_T(x)$ for $T$ equal to a specific, accumulated word library, not a claim about $L^*(x)$.

\section{Where the Two Deficits Diverge}
\label{app:counterexamples}

The Deficit Proxy Theorem's asymmetry is not merely an absence of proof; the converse direction can fail outright. The following toy construction gives an ontology with a large coding deficit and a near-zero distinguishability deficit, showing the case Section~\ref{sec:proxy} argues is left open is also, at least sometimes, actual.

\begin{example}[High coding deficit, low distinguishability deficit]
Let $X = \{x_1, \dots, x_n\}$ and let the true generating structure separate every pair: $D^*(x_i) = n - 1$ for each $i$. Let $T$ describe each $x_i$ by a fixed-width codeword of length $\lceil \log_2 n \rceil + m$ for some large padding constant $m > 0$ --- inefficient relative to $L^*(x_i)$, which needs only $\lceil \log_2 n \rceil$ bits to distinguish $n$ equally likely objects, so $\delta_T(x_i) \ge m$ for every $i$, growing without bound as $m$ grows. But the codewords remain pairwise distinct for every value of $m$: padding a codeword does not merge it with any other, so $T$ still separates every $x_i$ from every $x_j$, giving $D_T(x_i) = n - 1 = D^*(x_i)$ and therefore $\Delta_T(x_i) = 0$ regardless of how large $m$, and hence $\delta_T(x_i)$, is made.
\end{example}

The construction is deliberately trivial --- padding is not a realistic model of why real representations carry excess length --- but its triviality is the point: nothing in the definitions of the two deficits rules out this combination, so no additional theorem is needed to show the combination is possible, only an example showing it actually occurs. A more realistic version of the same divergence is the ordinary condition of a verbose but complete encoding: any description language that never merges distinct objects, however wastefully it encodes them, has $\Delta_T(x) = 0$ by construction, whatever $\delta_T(x)$ happens to be. Section~\ref{sec:forth}'s Forth library and Section~\ref{sec:ml}'s superposed neural representations are offered as the substantive, non-toy versions of this same structural possibility; this appendix's role is now to confirm that the possibility is real rather than merely unproven-against, with Proposition~\ref{prop:nonident} in Section~\ref{sec:nonident} superseding it as the general result: this example is the single point $k = |X/\!\sim|-1$, $\ell \to \infty$ on the surface that proposition characterizes in full.

\section{Consequences of Non-Identification}
\label{app:consequences}

The results of this paper admit several consequences stronger than the claim that compression is an unreliable heuristic for understanding. They concern the possibility of scalar theories of representation as such, and they follow from Proposition~\ref{prop:nonident} with no further assumptions beyond those already in use.

\begin{proposition}[No Universal Compression Ordering]
\label{prop:no-universal-order}
There is no nonconstant function $F : \mathbb{N} \to \mathbb{R}$ of description length alone such that, for all ontologies $T_1, T_2$ and all $q \in Q$, $F(L_{T_1}(q)) > F(L_{T_2}(q))$ implies $D_{T_1}(q) > D_{T_2}(q)$.
\end{proposition}

\begin{proof}
Suppose such an $F$ exists. Since $F$ is nonconstant, choose $\ell_1, \ell_2$ with $F(\ell_1) > F(\ell_2)$. By Proposition~\ref{prop:nonident}(i), construct $T_1$ with $L_{T_1}(q) = \ell_1$, $D_{T_1}(q) = 0$, and $T_2$ with $L_{T_2}(q) = \ell_2$, $D_{T_2}(q) = |Q|-1$; this choice of $(\ell,k)$ pairs is available regardless of the numerical relationship between $\ell_1$ and $\ell_2$, since part (i) fixes $\ell$ and $k$ independently. Then $F(L_{T_1}(q)) > F(L_{T_2}(q))$ holds by construction, while $D_{T_1}(q) = 0 < |Q|-1 = D_{T_2}(q)$ contradicts the assumed implication.
\end{proof}

No rescaling or nonlinear transform of description length repairs the proxy; the failure is not a bad choice of $F$, since Proposition~\ref{prop:no-universal-order} rules out every nonconstant $F$ at once.

\begin{corollary}[No Scalar Recovery from Description Length]
\label{cor:no-scalar-recovery}
There is no function $f : \mathbb{N} \to \mathbb{N}$ such that $D_T(q) = f(L_T(q))$ for every ontology $T$ and every $q \in Q$.
\end{corollary}

\begin{proof}
By Proposition~\ref{prop:nonident}(i), fix $\ell$ and choose $T_1, T_2$ with $L_{T_1}(q) = L_{T_2}(q) = \ell$ but $D_{T_1}(q) \ne D_{T_2}(q)$---available since $|Q| \ge 3$ gives at least two distinct achievable values of $k$ at any fixed $\ell$. If $D_T(q) = f(L_T(q))$ held universally, $f(\ell)$ could not take two different values at once, a contradiction.
\end{proof}

The next result treats the three coordinates of Definition~\ref{def:threecoord} together rather than $L_T$ and $D_T$ alone. Each coordinate carries its own natural order: $L_T(q) \in \mathbb{N}$ under $\le$; $\mathcal{P}_T$ under partition refinement, $\mathcal{P}_{T_1} \preceq \mathcal{P}_{T_2}$ meaning every block of $\mathcal{P}_{T_1}$ is a subset of some block of $\mathcal{P}_{T_2}$, so that $\mathcal{P}_{T_1} \preceq \mathcal{P}_{T_2}$ says $T_1$'s partition is at least as fine; and $\mathcal{O}_T$ under set inclusion $\subseteq$. Each of these is a standard partial order on its own, and their product is a partial order on the triple.

\begin{definition}[Representational Dominance]
\label{def:dominance}
For ontologies $T_1, T_2$ over the same $Q$, write $T_1 \succeq T_2$ when $L_{T_1}(q) \le L_{T_2}(q)$, $\mathcal{P}_{T_1} \preceq \mathcal{P}_{T_2}$, and $\mathcal{O}_{T_1} \supseteq \mathcal{O}_{T_2}$, with at least one inequality strict for strict dominance $T_1 \succ T_2$.
\end{definition}

\begin{proposition}[Incomparability of Representations]
\label{prop:incomparability}
Representational dominance is not a total order: there exist $T_1, T_2$ with neither $T_1 \succeq T_2$ nor $T_2 \succeq T_1$.
\end{proposition}

\begin{proof}
Take $T_1, T_2$ to be the ontologies of encoders $F_3$ and $F_5$ from Section~\ref{sec:nonident}'s worked example, both with $\mathcal{O}_{T_1} = \mathcal{O}_{T_2} = \emptyset$, so the operational coordinate is tied and the comparison reduces to $L_T$ and $\mathcal{P}_T$. Over $Q = \{a,b,c\}$, $F_3$ induces $\mathcal{P}_{F_3} = \{\{a,b\},\{c\}\}$ and $F_5$ induces $\mathcal{P}_{F_5} = \{\{a\},\{b\},\{c\}\}$, the discrete partition. $\mathcal{P}_{F_5} \preceq \mathcal{P}_{F_3}$: every singleton block of $\mathcal{P}_{F_5}$ is trivially a subset of some block of $\mathcal{P}_{F_3}$. But $\mathcal{P}_{F_3} \not\preceq \mathcal{P}_{F_5}$: the block $\{a,b\}$ of $\mathcal{P}_{F_3}$ is not a subset of any single block of the discrete $\mathcal{P}_{F_5}$. Now check both directions of dominance. $F_5 \succeq F_3$ requires $L_{F_5}(a) \le L_{F_3}(a)$, i.e.\ $4 \le 2$, which is false. $F_3 \succeq F_5$ requires $L_{F_3}(a) \le L_{F_5}(a)$ (true, $2 \le 4$) \emph{and} $\mathcal{P}_{F_3} \preceq \mathcal{P}_{F_5}$, which was just shown false. Neither direction holds.
\end{proof}

$F_3$ wins on description economy; $F_5$ wins on distinguishing structure; the comparison does not resolve without a further judgment about how the two are to be traded against each other. Better representations, in the sense this paper's apparatus measures, need not form a ranking.

\begin{definition}[Representational Pareto Frontier]
The representational Pareto frontier is $\mathfrak{F} = \{T : \nexists\, T' \text{ with } T' \succ T\}$. Members of $\mathfrak{F}$ cannot be ranked against one another without a preference specifying how description economy, distinguishing structure, and operational repertoire are to be weighed against one another; Proposition~\ref{prop:incomparability} shows $\mathfrak{F}$ is not generally a singleton or a chain.
\end{definition}

\begin{proposition}[Impossibility of Compression-Only Understanding Metrics]
\label{prop:no-understanding-metric}
Suppose a functional $U(T,q)$ is required to satisfy $D_{T_1}(q) > D_{T_2}(q) \implies U(T_1,q) > U(T_2,q)$ for all $T_1, T_2, q$. Then no function $g : \mathbb{N} \to \mathbb{R}$ satisfies $U(T,q) = g(L_T(q))$ for every ontology $T$.
\end{proposition}

\begin{proof}
By Proposition~\ref{prop:nonident}(i), fix $\ell$ and choose $T_1, T_2$ with $L_{T_1}(q) = L_{T_2}(q) = \ell$ and $D_{T_1}(q) > D_{T_2}(q)$. If $U(T,q) = g(L_T(q))$ held, then $U(T_1,q) = g(\ell) = U(T_2,q)$, contradicting the required $U(T_1,q) > U(T_2,q)$.
\end{proof}

This result does not depend on any particular definition of understanding; it says only that if a notion of understanding is required to be sensitive to distinguishing capacity at all, no version of it can be recovered as a function of description length alone. The temporal analogue---that compression progress does not identify distinguishing progress, and so cannot identify progress in any $U$ sensitive to $D_T$ in this sense---is Proposition~\ref{prop:dynamic}, already established in Section~\ref{sec:nonident}; it is not re-derived here.

\emph{Scalarization requires an external objective.} Once $L_T$, $\mathcal{P}_T$, and $\mathcal{O}_T$ are recognized as independently variable, any scalar ranking of representations takes the form $U(T) = F(L_T, \mathcal{P}_T, \mathcal{O}_T; \lambda)$ for some weighting $\lambda$ specifying how the three coordinates trade against one another. Nothing in $L_T$, $\mathcal{P}_T$, or $\mathcal{O}_T$ themselves determines $\lambda$; it is supplied from outside the representation, by whichever task, loss, or value judgment the comparison is actually being made for. This is an observation about what any scalarization must import, not a theorem with a proof of its own, and it is worth stating plainly for that reason: the mathematical content of this paper is that compression cannot supply a universal ordering of representations without such an import, not that compression is irrelevant to any particular one, chosen for a particular purpose. Representations, in the sense $\mathfrak{R} = \mathfrak{L} \times \mathfrak{P} \times \mathfrak{O}$ collects them---economy ordered oppositely to cost, partitions ordered by refinement, operations ordered by inclusion---occupy a product space, not a line, and the natural comparison relation on a product space is a product partial order, not a scalar one. In the unconstrained case, representation quality is not a scalar,
\[
T \; \mapsto \; \bigl(L_T, \mathcal{P}_T, \mathcal{O}_T\bigr),
\]
and comparison between representations is naturally partial rather than total. The pair $F_3, F_5$ constructed in Proposition~\ref{prop:incomparability} is not a hypothetical illustration of this fact but a finite witness to it, drawn from the paper's own apparatus: incomparability is not merely consistent with the definitions, it actually occurs inside the smallest example the paper builds.

The obstruction this appendix identifies is accordingly not that the correct scalar measure of representational quality has yet to be discovered. Under the assumptions of this paper, the representational space is not intrinsically scalar in the first place. Any map $U : \mathfrak{R} \to \mathbb{R}$ that totally orders otherwise incomparable representations necessarily adds a comparison rule not supplied by description length, partition refinement, or operational inclusion alone. This does not make such maps illegitimate: accuracy, minimum-description-length scores, computational cost, interpretability metrics, and task loss are all defensible scalar objectives for a given purpose. What changes is their status. Each is a projection of $\mathfrak{R}$ chosen for a reason external to $\mathfrak{R}$ itself, not a discovery of the one-dimensional quantity ``representational quality'' that description length was, on the assumption this paper has been examining, presumed to already be measuring.

\begin{center}
\fbox{\parbox{0.85\linewidth}{\centering measurement requires projection; projection does not imply identification}}
\end{center}

``Shorter'' is one coordinate of $\mathfrak{R}$. This paper has not shown that ``better understood'' is another; it has shown that no argument considered here, or ordinarily offered elsewhere, establishes that the two coordinates coincide.

\newpage
\section*{References}
\addcontentsline{toc}{section}{References}

{\renewcommand{\refname}{}
\begin{thebibliography}{99}

\bibitem{distgeom}
Flyxion.
\newblock 2026.
\newblock \emph{Distinguishability Geometry: Projection, Embedding, Revision, and Transport}.
\newblock \url{https://standardgalactic.github.io/playfloor/working/distinguishability_geometry.pdf}.

\bibitem{compexp}
Flyxion.
\newblock 2026b.
\newblock \emph{Compression After Expansion: On Skill, Misattribution, and the Geometry of Work}.
\newblock \url{https://standardgalactic.github.io/alphabet/screenplay/compression_after_expansion.pdf}.

\bibitem{ziembik2025}
Ziembik, Krzysztof A.
\newblock 2025.
\newblock ``Forth: The Best Programming Language for the End of the World? A Case Study on A.D.\ 2044 (1991).''
\newblock \emph{Proceedings of the 36th ACM Conference on Hypertext and Social Media}.

\bibitem{rissanen1978}
Rissanen, Jorma.
\newblock 1978.
\newblock ``Modeling by Shortest Data Description.''
\newblock \emph{Automatica} 14(5): 465--471.

\bibitem{schmidhuber2008}
Schmidhuber, J\"urgen.
\newblock 2008.
\newblock ``Driven by Compression Progress: A Simple Principle Explains Essential Aspects of Subjective Beauty, Novelty, Surprise, Interestingness, Attention, Curiosity, Creativity, Art, Science, Music, Jokes.''
\newblock \emph{arXiv preprint} arXiv:0812.4360.
\newblock Short version published as ``Simple Algorithmic Theory of Subjective Beauty, Novelty, Surprise, Interestingness, Attention, Curiosity, Creativity, Art, Science, Music, Jokes,'' \emph{Journal of SICE} 48(1): 21--32, 2009.

\bibitem{higgins2017}
Higgins, Irina, Lo\"ic Matthey, Arka Pal, Christopher Burgess, Xavier Glorot, Matthew Botvinick, Shakir Mohamed, and Alexander Lerchner.
\newblock 2017.
\newblock ``beta-VAE: Learning Basic Visual Concepts with a Constrained Variational Framework.''
\newblock \emph{International Conference on Learning Representations (ICLR)}.

\bibitem{bricken2023}
Bricken, Trenton, Adly Templeton, Joshua Batson, Brian Chen, Adam Jermyn, Tom Conerly, Nicholas L.\ Turner, et al.
\newblock 2023.
\newblock ``Towards Monosemanticity: Decomposing Language Models with Dictionary Learning.''
\newblock \emph{Transformer Circuits Thread}.

\bibitem{gopnik2020}
Gopnik, Alison.
\newblock 2020.
\newblock ``Childhood as a Solution to Explore--Exploit Tensions.''
\newblock \emph{Philosophical Transactions of the Royal Society B: Biological Sciences} 375(1803): 20190502.

\bibitem{silva2026merits}
Silva, David William.
\newblock 2026.
\newblock ``Merits of Geometric Algebra Applied to Cryptography and Machine Learning.''
\newblock \emph{Philosophical Transactions of the Royal Society A} 384: 20250107.
\newblock \url{https://doi.org/10.1098/rsta.2025.0107}.

\bibitem{sotudeh2022tstr}
Sotudeh, Sajad, and Nazli Goharian.
\newblock 2022.
\newblock ``TSTR: Too Short to Represent, Summarize with Details! Intro-Guided Extended Summary Generation.''
\newblock \emph{arXiv preprint} arXiv:2206.00847.

\bibitem{forthsdilemma}
Anonymous.
\newblock c.~2000.
\newblock ``The Forth's Dilemma, When Great Firms Cause New Technologies to Fail.''
\newblock Archived at the Internet Archive Wayback Machine, capture of \url{http://www.msmisp.com/futuretest/Forth's_Dilemma.htm}.

\end{thebibliography}
}

\end{document}

